%! Algebra 1 Priority Standards Quiz --- Version B (blank)
% GENERATED BODY: the blank and key bodies are byte-identical except the header title;
% answers live in \slot / \optok / work, which the preamble shows or hides.
% Standards: 2023 VA SOL Algebra 1 — A.EO.1, A.EO.2, A.EO.3, A.EO.4, A.EI.1, A.EI.2, A.EI.3,
% A.F.1, A.F.2. A.ST.1 (regression, scatterplots) is deliberately left out.
\documentclass[12pt]{article}
\usepackage{algebra2-article}
\usepackage{algebra2-boxes}
\usepackage{grading}
\gradingassignment{alg2-u01-algebra1-quiz-b}
\gradingcourse{algebra_2}
\usepackage{algebra2-graded}
% Coordinate grid: {xmin}{xmax}{ymin}{ymax}
\newcommand{\lgrid}[4]{%
  \draw[step=1, linegray!40, very thin] (#1,#3) grid (#2,#4);
  \draw[->, semithick, charcoal] (#1,0)--(#2,0) node[below right, font=\scriptsize]{$x$};
  \draw[->, semithick, charcoal] (0,#3)--(0,#4) node[left, font=\scriptsize]{$y$};}
% Part-divider strip (headlinebox takes one arg: the background color)
\newcommand{\parthead}[1]{%
  \boxguard[10]\vspace{6pt}\noindent
  \begin{headlinebox}{forest}\color{white}\bfseries #1\end{headlinebox}%
  \vspace{4pt}}
% THE WORK RULE for a test: \slot{width}{answer} and \opt / \optok are the ONLY
% lines that differ between blank and key, and they differ in the preamble, not the body.
\newcommand{\slot}[2]{\underline{\makebox[#1]{\vphantom{\textbf{#2}}}}}% the key's answer height, so blank and key keep the same lines
\newcommand{\opt}[2]{(#1)~#2}
\newcommand{\optok}[2]{(#1)~#2}

\begin{document}

\pageheader{Algebra 1 Review}{Priority Standards Quiz --- Version B}

\vspace{0.06in}
\noindent\textbf{Instructions:} No calculator. Show your work where space is given. Circle the
letter for multiple choice. Leave radicals in simplest form.
\hfill\textbf{25 questions, 4 pts each = 100 pts}

% ======================================================================
\parthead{Part A --- Expressions, Exponents \& Radicals \hfill\normalfont\small A.EO.1, A.EO.3, A.EO.4}

\begin{enumerate}[leftmargin=1.9em, itemsep=10pt, label=\arabic*.]
  \item Evaluate $3x^2-|x+6|$ when $x=-4$.
    \begin{work}
      3(-4)^2-|-4+6| &= 3(16)-|2| \\
      &= 48-2
    \end{work}
    \hfill Value: \inblank{\slot{2.4cm}{$46$}}{\gradedblank[id=q1, type=short, points=4, answer={46}, check=numeric]{2.4cm-2pt}{13.5pt}}

  \item \letterbox{\gradedblank[id=q2, type=choice, points=4, answer={C}]{0.4in}{13.5pt}}Which expression represents ``seven more than twice a number $n$''?
    \\*[4pt]
    \opt{A}{$2+7n$} \hfill \opt{B}{$2(n+7)$} \hfill \optok{C}{$2n+7$} \hfill \opt{D}{$7-2n$}

  \item Simplify $(3x^2y^3)^2\cdot y^{-5}$. Use only positive exponents.
    \begin{work}
      (3x^2y^3)^2\cdot y^{-5} &= 9x^4y^6\cdot y^{-5}
    \end{work}
    \hfill Answer: \inblank{\slot{3.0cm}{$9x^4y$}}{\gradedblank[id=q3, type=short, points=4, answer={9x^4y}, check=expression]{3.0cm-2pt}{13.5pt}}

  \item Simplify $\dfrac{18m^3n^4}{6m^7n}$. Use only positive exponents.
    \begin{work}
      \frac{18}{6}\cdot m^{3-7}\cdot n^{4-1} &= 3m^{-4}n^3
    \end{work}
    \hfill Answer: \inblank{\slot{3.0cm}{$\dfrac{3n^3}{m^4}$}}{\gradedblank[id=q4, type=short, points=4, answer={3n^3/m^4}, check=expression, display={$\dfrac{3n^3}{m^4}$}]{3.0cm-2pt}{13.5pt}}

  \item Simplify $\sqrt{75}$.
    \begin{work}
      \sqrt{75} &= \sqrt{25}\cdot\sqrt{3}
    \end{work}
    \hfill Answer: \inblank{\slot{3.0cm}{$5\sqrt{3}$}}{\gradedblank[id=q5, type=short, points=4, answer={5sqrt(3)}, check=expression, display={$5\sqrt{3}$}]{3.0cm-2pt}{13.5pt}}

  \item \letterbox{\gradedblank[id=q6, type=choice, points=4, answer={D}]{0.4in}{13.5pt}}Which is equivalent to $3\sqrt{2}+4\sqrt{8}$?
    \\*[4pt]
    \opt{A}{$7\sqrt{10}$} \hfill \opt{B}{$7\sqrt{2}$} \hfill \opt{C}{$19\sqrt{2}$} \hfill
    \optok{D}{$11\sqrt{2}$}
\end{enumerate}

% ======================================================================
\parthead{Part B --- Polynomials \& Factoring \hfill\normalfont\small A.EO.2}

\begin{enumerate}[leftmargin=1.9em, itemsep=10pt, label=\arabic*., start=7]
  \item Subtract: $(4x^2+3x-5)-(2x^2-x+3)$.
    \begin{work}
      &= 4x^2+3x-5-2x^2+x-3
    \end{work}
    \hfill Answer: \inblank{\slot{3.6cm}{$2x^2+4x-8$}}{\gradedblank[id=q7, type=short, points=4, answer={2x^2+4x-8}, check=expression]{3.6cm-2pt}{13.5pt}}

  \item Multiply: $(3x+2)(x-5)$.
    \begin{work}
      &= 3x^2-15x+2x-10
    \end{work}
    \hfill Answer: \inblank{\slot{3.6cm}{$3x^2-13x-10$}}{\gradedblank[id=q8, type=short, points=4, answer={3x^2-13x-10}, check=expression]{3.6cm-2pt}{13.5pt}}

  \item Factor: $x^2+3x-18$.
    \\[1.4cm]
    \hspace*{\fill}Answer: \inblank{\slot{3.6cm}{$(x+6)(x-3)$}}{\gradedblank[id=q9, type=short, points=4, answer={(x+6)(x-3)}, check=expression]{3.6cm-2pt}{13.5pt}}

  \item Factor: $3x^2+10x+8$.
    \\[1.4cm]
    \hspace*{\fill}Answer: \inblank{\slot{3.6cm}{$(3x+4)(x+2)$}}{\gradedblank[id=q10, type=short, points=4, answer={(3x+4)(x+2)}, check=expression]{3.6cm-2pt}{13.5pt}}

  \item \letterbox{\gradedblank[id=q11, type=choice, points=4, answer={B}]{0.4in}{13.5pt}}Which is the \textbf{complete} factorization of $2x^2-32$?
    \\*[4pt]
    \opt{A}{$2(x^2-16)$} \hfill \optok{B}{$2(x-4)(x+4)$} \hfill \opt{C}{$2(x-4)^2$} \hfill
    \opt{D}{$(2x-8)(x+4)$}
\end{enumerate}

% ======================================================================
\parthead{Part C --- Linear Equations, Inequalities \& Systems \hfill\normalfont\small A.EI.1, A.EI.2}

\begin{enumerate}[leftmargin=1.9em, itemsep=10pt, label=\arabic*., start=12]
  \item Solve: $5(x-1)=2x+13$.
    \begin{work}
      5x-5 &= 2x+13 \\
      3x &= 18
    \end{work}
    \hfill Solution: \inblank{\slot{2.4cm}{$x=6$}}{\gradedblank[id=q12, type=short, points=4, answer={x=6}, check=expression]{2.4cm-2pt}{13.5pt}}

  \item \letterbox{\gradedblank[id=q13, type=choice, points=4, answer={A}]{0.4in}{13.5pt}}Which is the solution of $-2x+7>15$?
    \\*[4pt]
    \optok{A}{$x<-4$} \hfill \opt{B}{$x>-4$} \hfill \opt{C}{$x<4$} \hfill
    \opt{D}{$x>4$}

  \item Solve $P=2l+2w$ for $w$.
    \begin{work}
      P-2l &= 2w
    \end{work}
    \hfill $w=$ \inblank{\slot{2.4cm}{$\dfrac{P-2l}{2}$}}{\gradedblank[id=q14, type=short, points=4, answer={(P-2l)/2}, check=expression, display={$\dfrac{P-2l}{2}$}]{2.4cm-2pt}{13.5pt}}

  \item \letterbox{\gradedblank[id=q15, type=choice, points=4, answer={D}]{0.4in}{13.5pt}}How many solutions does $3(x-2)=3x-6$ have?
    \\*[4pt]
    \opt{A}{exactly one} \hfill \opt{B}{no solution} \hfill \opt{C}{exactly two} \hfill
    \optok{D}{infinitely many}

  \item Solve the system. \quad
    $\begin{cases} y=x+3 \\ 2x+y=12 \end{cases}$
    \begin{work}
      2x+(x+3) &= 12 \\
      3x &= 9 \quad\Rightarrow\quad x=3 \\
      y &= 3+3 = 6
    \end{work}
    \hfill Solution: \inblank{\slot{2.8cm}{$(3,\,6)$}}{\gradedblank[id=q16, type=short, points=4, answer={(3,6)}, check=text, display={$(3,\,6)$}]{2.8cm-2pt}{13.5pt}}

  \item \letterbox{\gradedblank[id=q17, type=choice, points=4, answer={B}]{0.4in}{13.5pt}}Which ordered pair is a solution of $y<2x-1$?
    \\*[4pt]
    \opt{A}{$(0,\,0)$} \hfill \optok{B}{$(2,\,1)$} \hfill \opt{C}{$(1,\,3)$} \hfill
    \opt{D}{$(-1,\,-2)$}
\end{enumerate}

% ======================================================================
\parthead{Part D --- Linear Functions \hfill\normalfont\small A.F.1}

\begin{enumerate}[leftmargin=1.9em, itemsep=10pt, label=\arabic*., start=18]
  \item For the line $4x+3y=12$, find the slope and the $y$-intercept.
    \begin{work}
      3y &= -4x+12 \\
      y &= -\frac{4}{3}x+4
    \end{work}
    \hfill Slope: \inblank{\slot{2.0cm}{$-\tfrac{4}{3}$}}{\gradedblank[id=q18a, type=short, points=2, answer={-4/3}, check=numeric, display={$-\tfrac{4}{3}$}]{2.0cm-2pt}{13.5pt}} \qquad $y$-intercept: \inblank{\slot{2.0cm}{$(0,\,4)$}}{\gradedblank[id=q18b, type=short, points=2, answer={(0,4)}, check=text, display={$(0,\,4)$}]{2.0cm-2pt}{13.5pt}}

  \item Write the equation, in $y=mx+b$ form, of the line \textbf{parallel} to $y=3x-2$ that
    passes through $(2,\,1)$.
    \begin{work}
      y-1 &= 3(x-2) \\
      y-1 &= 3x-6
    \end{work}
    \hfill Equation: \inblank{\slot{3.6cm}{$y=3x-5$}}{\gradedblank[id=q19, type=short, points=4, answer={y=3x-5}, check=expression]{3.6cm-2pt}{13.5pt}}

  \boxguard[9]
  \item \letterbox{\gradedblank[id=q20, type=choice, points=4, answer={A}]{0.4in}{13.5pt}}Which equation matches the graph?
    \\*[4pt]
    \begin{minipage}[c]{0.36\linewidth}\centering
    \begin{tikzpicture}[scale=0.36]
      \lgrid{-4}{4}{-4}{4}
      \draw[very thick, navy] (-1.333,-4)--(4,4);
      \fill[navy] (0,-2) circle (4pt); \fill[navy] (2,1) circle (4pt);
    \end{tikzpicture}
    \end{minipage}%
    \begin{minipage}[c]{0.6\linewidth}
    \optok{A}{$y=\tfrac{3}{2}x-2$} \\[6pt]
    \opt{B}{$y=\tfrac{2}{3}x-2$} \\[6pt]
    \opt{C}{$y=\tfrac{3}{2}x+2$} \\[6pt]
    \opt{D}{$y=-\tfrac{3}{2}x-2$}
    \end{minipage}

  \item For $f(x)=4x-8$, find $f(-3)$, and find the value of $x$ for which $f(x)=0$.
    \begin{work}
      f(-3) &= 4(-3)-8 = -20 \\
      0 &= 4x-8 \quad\Rightarrow\quad x=2
    \end{work}
    \hfill $f(-3)=$ \inblank{\slot{1.6cm}{$-20$}}{\gradedblank[id=q21a, type=short, points=2, answer={-20}, check=numeric]{1.6cm-2pt}{13.5pt}} \qquad $x=$ \inblank{\slot{1.6cm}{$2$}}{\gradedblank[id=q21b, type=short, points=2, answer={2}, check=numeric]{1.6cm-2pt}{13.5pt}}
\end{enumerate}

% ======================================================================
\parthead{Part E --- Functions \& Quadratics \hfill\normalfont\small A.F.2, A.EI.3}

\begin{enumerate}[leftmargin=1.9em, itemsep=10pt, label=\arabic*., start=22]
  \item \letterbox{\gradedblank[id=q22, type=choice, points=4, answer={C}]{0.4in}{13.5pt}}Which relation is \textbf{not} a function?
    \\*[4pt]
    \opt{A}{$\{(1,3),\,(2,4),\,(3,5)\}$} \hfill \opt{B}{$\{(4,1),\,(5,1),\,(6,1)\}$}
    \\[6pt]
    \optok{C}{$\{(3,1),\,(3,2),\,(4,5)\}$} \hfill \opt{D}{$\{(0,0),\,(1,-1),\,(2,-2)\}$}

  \boxguard[12]
  \item The graph of $f(x)=x^2+4x+3$ is shown. Read its key features.
    \\[4pt]
    \begin{minipage}[c]{0.40\linewidth}\centering
    \begin{tikzpicture}[scale=0.36]
      \lgrid{-5}{2}{-2}{5}
      \draw[very thick, navy, domain=-4.4:0.4, smooth, samples=40] plot (\x, {\x*\x+4*\x+3});
      \foreach \p in {(-3,0),(-1,0),(-2,-1),(0,3)} \fill[navy] \p circle (4pt);
    \end{tikzpicture}
    \end{minipage}%
    \begin{minipage}[c]{0.58\linewidth}
    Zeros: \inblank{\slot{3.4cm}{$x=-3$, \ $x=-1$}}{\gradedblank[id=q23a, type=short, points=1, answer={-3, -1}, check=set, display={$x=-3$, \ $x=-1$}]{3.4cm-2pt}{13.5pt}} \\[12pt]
    Vertex: \inblank{\slot{2.4cm}{$(-2,\,-1)$}}{\gradedblank[id=q23b, type=short, points=1, answer={(-2,-1)}, check=text, display={$(-2,\,-1)$}]{2.4cm-2pt}{13.5pt}} \\[12pt]
    $y$-intercept: \inblank{\slot{2.4cm}{$(0,\,3)$}}{\gradedblank[id=q23c, type=short, points=1, answer={(0,3)}, check=text, display={$(0,\,3)$}]{2.4cm-2pt}{13.5pt}} \\[12pt]
    Range: \inblank{\slot{2.4cm}{$y\ge -1$}}{\gradedblank[id=q23d, type=short, points=1, answer={y>=-1}, check=expression, display={$y\ge -1$}]{2.4cm-2pt}{13.5pt}}
    \end{minipage}

  \item Solve by factoring: $x^2-x-12=0$.
    \begin{work}
      (x-4)(x+3) &= 0 \\
      x-4=0 \quad&\text{or}\quad x+3=0
    \end{work}
    \hfill Solutions: \inblank{\slot{3.6cm}{$x=4$ or $x=-3$}}{\gradedblank[id=q24, type=short, points=4, answer={4, -3}, check=set, display={$x=4$ or $x=-3$}]{3.6cm-2pt}{13.5pt}}

  \item Solve: $x^2+6x+2=0$. Give exact answers in simplest radical form.
    \begin{work}
      x &= \frac{-6\pm\sqrt{6^2-4(1)(2)}}{2(1)} = \frac{-6\pm\sqrt{28}}{2} \\
      &= \frac{-6\pm 2\sqrt{7}}{2}
    \end{work}
    \hfill Solutions: \inblank{\slot{3.6cm}{$x=-3\pm\sqrt{7}$}}{\gradedblank[id=q25, type=short, points=4, answer={-3+sqrt(7), -3-sqrt(7)}, check=set, display={$x=-3\pm\sqrt{7}$}]{3.6cm-2pt}{13.5pt}}
\end{enumerate}

\end{document}
